\documentclass[11pt]{article}
\usepackage{mathblatt}
\makeatletter
\setlength{\mbpfbe}{0mm}% keine Punkte (Bauregel 6.4)
\makeatother
% Rechenraum als Karo (Bauregel 6.9): n Rechenschritte = n mal 10 mm
\newcommand{\karo}[1]{\par\smallskip\noindent\begin{tikzpicture}
  \clip (0,0) rectangle ({\linewidth-1mm},{#1*10mm});
  \draw[black!16,line width=0.3pt,step=5mm] (0,0) grid ({\linewidth},{#1*10mm});
  \end{tikzpicture}\par}
\newcommand{\linie}[1]{\,{\color{black!45}\rule[-1pt]{#1}{0.4pt}}\,}
\newcommand{\pkt}{\textperiodcentered{}}
\begin{document}

\pfheftstil
\pflaufkopf[P10 \pkt{} mehr]{Diagramme \pkt{} Lösungen}{PS40-3}\pfseitevon
\pfheftkopf{Diagramme}{Lösungen \pkt{} mehr}
\begin{pfloesungb}{}
\lzb{1.}{}{Flüssigkeiten 13 \%}{im Uhrzeigersinn ab 12 Uhr: Getreide, Flüssigkeiten, Schrott, Sonstiges, Container; Flüssigkeiten 13 \%}
\end{pfloesungb}
\begin{pfloesungb}{}
\lzb{2.}{}{Streifen mit Abschnitten 0,6 cm (nie), 1,8 cm (1–3-mal), 9,0 cm (4–6-mal), 3,6 cm (täglich)}{Streifen mit Abschnitten 0,6 cm (nie), 1,8 cm (1–3-mal), 9,0 cm (4–6-mal), 3,6 cm (täglich)}
\end{pfloesungb}
\begin{pfloesungb}{}
\lzb{3.}{}{waagerechter Balken bei Do bis 155}{waagerechter Balken bei Do bis 155}
\end{pfloesungb}
\begin{pfloesungb}{}
\lzb{4.}{}{grauer Sektor: viertes oder weiteres Kind (4,4 \% von 360° $\approx$ 16°); drittes Kind: 11,8 \% von 360° $\approx$ 42°}{grauer Sektor: viertes oder weiteres Kind (4,4 \% von 360° $\approx$ 16°); drittes Kind: 11,8 \% von 360° $\approx$ 42°}
\end{pfloesungb}
\begin{pfloesungb}{}
\lzb{5.}{}{Fisch 150 : 20 = 7,5 Kästchen}{Achse 0, 20, 40, … (20 je Gitterlinie); E und D: Kartoffelgerichte und Gemüsegerichte (Reihenfolge beliebig), C: Pizza, B: Schnitzel, A: Fischgerichte; Fisch 150 : 20 = 7,5 Kästchen}
\end{pfloesungb}
\begin{pfloesungb}{}
\lzb{6.}{}{Aussage 1 falsch (60 \% sind mehr als die Hälfte); Aussage 2 falsch (jede 12. wäre etwa 8,3 \%, tatsächlich 12 \%)}{Aussage 1 falsch (60 \% sind mehr als die Hälfte); Aussage 2 falsch (jede 12. wäre etwa 8,3 \%, tatsächlich 12 \%)}
\end{pfloesungb}
\begin{pfloesungb}{}
\lzb{7.}{}{z. B.: „extrem abnimmt“ ist falsch, die abgeschnittene Achse übertreibt – Rückgang von rund 47 \% auf 40 \% in 4 Jahren, also nur 7 Prozentpunkte; oder: „in wenigen Jahren kein Jugendlicher mehr“ ist falsch – 2016 lesen noch 40 \%, bei gleichem Tempo erst in rund 23 Jahren}{z. B.: „extrem abnimmt“ ist falsch, die abgeschnittene Achse übertreibt – Rückgang von rund 47 \% auf 40 \% in 4 Jahren, also nur 7 Prozentpunkte; oder: „in wenigen Jahren kein Jugendlicher mehr“ ist falsch – 2016 lesen noch 40 \%, bei gleichem Tempo erst in rund 23 Jahren}
\end{pfloesungb}
\begin{pfloesungb}{}
\lzb{8.}{}{echter Zuwachs je Jahr nur 4 €}{abgeschnittene, gedehnte y-Achse (beginnt bei 1040 € statt bei 0): Unterschiede von wenigen Euro erscheinen als Verdreifachung; echter Zuwachs je Jahr nur 4 €}
\end{pfloesungb}
\end{document}